% Quarantine-only source-aware proof.
% Date: 29 August 2026.
% This file is not a live-book input and makes no novelty or priority claim.

\begin{proof}[Proof of
\cref{lem:p21p3-componentwise-normal-rigid-zero-atom}]
We first reduce normality to the precise local coordinate needed in the
estimate.  Since $X$ is good, $\mathcal O_{X,y}$ is a noetherian henselian
excellent local ring by \cite[\S11.1.10]{CLD25}.  A rigid point is Zariski
closed in a sufficiently small strict affinoid neighborhood.  Indeed, the
fraction field of the quotient by the evaluation kernel embeds densely in
the finite-dimensional $K$-vector space $\mathscr H(y)$.  It is therefore
finite-dimensional over $K$, as is its $K$-subalgebra given by the affinoid
quotient.  That quotient is a finite-dimensional domain over $K$, hence a
field, so the evaluation kernel is maximal.  The central dimension of the
germ is consequently zero, and the local dimension formula
\cite[(11.5.5.1)]{CLD25}, together with the assumed pure local dimension,
gives
\[
                  \dim_{\mathrm{Krull}}\mathcal O_{X,y}=1.
\]
The normal local ring $\mathcal O_{X,y}$ is a domain.  The equivalence
between a one-dimensional noetherian normal local domain and a discrete
valuation ring is
\cite[\href{https://stacks.math.columbia.edu/tag/00PD}{Tag~00PD}]
{stacks-project}.  We may therefore choose a uniformizer
$\pi\in\mathcal O_{X,y}$.

Put $X^\circ:=\operatorname{Int}(X/K)$ and
$\Omega^\circ:=\Omega\cap X^\circ$.  These sets are open, and the
boundary-composition identity \cite[(11.4.1.3)]{CLD25} shows that every open
analytic subspace of $X^\circ$ is boundaryless.  After shrinking around
$y$, we may consequently work in a boundaryless good separated strict open
set
\[
                         U\subseteq\Omega^\circ
\]
on which $X$ is purely one-dimensional and normal and on which $\pi$ is
represented by an analytic function $f$ satisfying $f(y)=0$.  Such a normal
shrinking exists because the ordinary normal locus is open by
\cite[\S11.3.1]{CLD25}.  The germ of the fibre of
$f:U\to\mathbf A_K^{1,\mathrm{an}}$ over $0$ is zero-dimensional at $y$,
because
\[
                  \mathcal O_{X,y}/(f)
                  =\mathcal O_{X,y}/(\pi)
\]
is a field.  Formula (11.4.1.3) also gives
$y\in\operatorname{Int}(U/\mathbf A_K^{1,\mathrm{an}})$.  Thus
\cite[Lemma~11.6.5(a)]{CLD25} makes $f$ finite at $y$.

We spell out the flatness point, since ordinary module-flatness must not be
silently substituted for analytic flatness.  Let
\[
       A:=\mathcal O_{\mathbf A_K^{1,\mathrm{an}},0},
       \qquad B:=\mathcal O_{X,y}.
\]
The ring $A$ is a discrete valuation ring with uniformizer $T$, and the
local homomorphism $A\to B$ sends $T$ to $\pi$.  It is injective: if
$0\ne a\in A$, write $a=T^m v$ with $v$ a unit; then
$a(f)=\pi^m v(f)\ne0$ in the domain $B$.  Since $f$ is finite at $y$, the
$A$-module $B$ is finite by \cite[\S11.4.4]{CLD25}.  It is torsion-free,
hence finite free over the discrete valuation ring $A$.  Thus $B$ is flat
over $A$.  Since $y\in\operatorname{Int}(U/\mathbf A_K^{1,\mathrm{an}})$,
the local-ring criterion of \cite[Remark~11.4.7]{CLD25} upgrades this
module-flatness to analytic flatness of $f$ at $y$.  The local finite-flat
construction of \cite[\S11.6.3]{CLD25} now gives compact analytic
neighborhoods on which $f$ is finite and flat and on which $y$ is the only
point over $0$.  Restricting the target to a sufficiently small closed disc
of admissible radius $R_0>0$, and the source to its inverse image, therefore
produces a compact analytic neighborhood $Q_{R_0}\subset U$ of $y$ and a
finite flat map of constant degree $d>0$
\begin{equation}
\label{eq:p213-normal-rigid-finite-tube}
          f:Q_{R_0}\longrightarrow D(0,R_0),
          \qquad f^{-1}(0)=\{y\}.
\end{equation}
A finite morphism has empty relative boundary.  By
\cite[(11.4.1.3) and Example~11.4.2]{CLD25}, the boundary of $Q_{R_0}$ is
therefore contained in the inverse image of the outer Gauss point
$\{|T|=R_0\}$.

Choose an admissible radius $0<R<R_0$ and a number
\[
             0<\varepsilon_0<\tfrac12\log(R_0/R).
\]
For every admissible $0<r<R$ such that
$\log(R/r)>4\varepsilon_0$, set
\[
                    s_R=-\log R,
                    \qquad s_r=-\log r.
\]
Choose a function $\rho_r\in C^\infty(\mathbf R,[0,1])$ which vanishes on
\[
 (-\infty,s_R-\varepsilon_0]\cup
 [s_r+\varepsilon_0,\infty),
\]
equals $1$ on $[s_R,s_r]$, is nondecreasing on
$[s_R-\varepsilon_0,s_R]$, and is nonincreasing on
$[s_r,s_r+\varepsilon_0]$.  Define
\[
 q_r(t):=\int_{-\infty}^{t}\rho_r(\tau)\,\mathrm d\tau,
 \qquad
 H_r:=\int_{\mathbf R}\rho_r(\tau)\,\mathrm d\tau.
\]
Then
\begin{equation}
\label{eq:p213-normal-rigid-radial-data}
 H_r\geq\log(R/r),
 \qquad
 \int_{\mathbf R}|q_r''(t)|\,\mathrm dt=2.
\end{equation}

On $Q_{R_0}\setminus\{y\}$ put
\[
                         \chi_r=q_r(-\log|f|).
\]
The function $q_r$ is zero on an outer collar of $Q_{R_0}$ and is constant
equal to $H_r$ for sufficiently large arguments.  We may therefore extend
$\chi_r$ by $H_r$ across $y$ and by zero to $U\setminus Q_{R_0}$.  The
extensions are made where the function is constant; on the remaining locus
$f$ is invertible.  The moment description of smooth functions in
\cite[Corollary~15.1.2]{CLD25} therefore shows that $\chi_r$ is a
nonnegative compactly supported CLD-smooth function on $U$.  Moreover,
\begin{equation}
\label{eq:p213-normal-rigid-value-at-y}
                         \chi_r(y)=H_r.
\end{equation}

Let $s=-\log|f|$.  On the two transition tubes one has
\begin{equation}
\label{eq:p213-normal-rigid-radial-dd}
       \mathrm d'\mathrm d''\chi_r
       =q_r''(s)\,\mathrm d's\wedge\mathrm d''s.
\end{equation}
The two tubes are disjoint, and $q_r''$ has a fixed sign on each one.
Apply the calibrated pushforward formula
\cite[Proposition~14.4.2]{CLD25} to the restriction of
\eqref{eq:p213-normal-rigid-finite-tube} over each corresponding target
annulus.  The identification of calibrated degree with finite-flat degree
in \cite[Lemma~13.3.8]{CLD25} shows that the pushforward factor is $d$;
this includes all local and residue-degree multiplicities.  Consequently
\begin{align}
 \left\|\mathrm d'\mathrm d''\chi_r\wedge\delta_U
 \right\|_{\mathrm{TV}}
 &=d\int_{\mathbf R}|q_r''(t)|\,\mathrm dt \notag\\
 &=2d.
\label{eq:p213-normal-rigid-radial-tv}
\end{align}
Neither \cite[Proposition~14.4.2]{CLD25} nor
\cite[Lemma~13.3.8]{CLD25} requires the source to be rig-smooth.

Because $\Omega^\circ$ is boundaryless, its integration current is closed
by \cite[\S7.7.4]{CLD25}.  Applying \cite[Lemma~8.5.5]{CLD25} with
$T=\delta_{\Omega^\circ}$ shows that the current
\[
                         \mu_u
       =\mathrm d'\mathrm d''[u|_{\Omega^\circ}]
\]
is positive on $\Omega^\circ$; as a positive current of bidimension
$(0,0)$ it is a Radon measure by \cite[Proposition~7.2.8]{CLD25}.  Its
restriction to the boundaryless open set $U$ is
\[
                 \mu_u|_U=\mathrm d'\mathrm d''[u|_U].
\]
Equivalently, one may apply the same two results with $T=\delta_U$
directly.  Positivity,
\eqref{eq:p213-normal-rigid-radial-data},
\eqref{eq:p213-normal-rigid-value-at-y}, the transposition identity
\cite[Lemma~7.7.5]{CLD25}, and
\eqref{eq:p213-normal-rigid-radial-tv} yield
\begin{align}
 \log(R/r)\,\mu_u(\{y\})
 &\leq H_r\mu_u(\{y\})                                      \notag\\
 &\leq\int_U\chi_r\,\mathrm d\mu_u                          \notag\\
 &=\bigl\langle[u|_U],\mathrm d'\mathrm d''\chi_r\bigr\rangle \notag\\
 &\leq 2d\sup_{Q_{R_0}}|u|.
\label{eq:p213-normal-rigid-final-bound}
\end{align}
The last line follows by first taking the absolute value of the pairing and
then using its total-variation bound.  Its right-hand side is independent of
$r$.  Since the valuation of $K$ is nontrivial, admissible positive radii
tend to zero.  Letting $r\to0$ in
\eqref{eq:p213-normal-rigid-final-bound} proves
\eqref{eq:p213-normal-rigid-zero-atom}.
\end{proof}
